\SourceLeaf{376}%
Zwecke, durch französische Ingenieurs ausgeführt. Um das Jahr 1811 war ein Dreiecksnetz von Fiume bis Turin vollendet, welches mithin in der Richtung eines Parallelkreises des 45sten Breitengrades sich über sieben Längengrade erstreckte. Um diese Arbeiten auch in höherer wissenschaftlicher Beziehung für die Kenntniss der Gestalt der Erde nützlich zu machen, beschloss das damalige französische Gouvernement, auf \textsc{Laplace}’s Antrag, dieses Dreiecksnetz im Westen bis zum atlantischen Meere erweitern und die zu einer Längengradmessung erforderlichen Operationen damit verbinden zu lassen. Die sofort mit Eifer angefangene, nachher durch die Zeitereignisse eine Zeitlang unterbrochene, bald aber wieder mit gleicher Thätigkeit fortgesetzte Arbeit war im J. 1818 so weit gediehen, dass das Dreiecksnetz über das französische Gebiet vom atlantischen Meere bei Bordeaux bis an die Grenze von Savoyen gemessen war. Es fehlte also, zur Vollendung des geodaetischen Theils, nur noch das in den Staaten des Königs von Sardinien liegende Stück. Das dortige und das Oesterreichische Gouvernement, beide die wissenschaftliche Wichtigkeit dieser grossartigen Unternehmung lebhaft anerkennend, beschlossen, durch eine aus Astronomen und Officieren beider Staaten zusammengesetzte Commission sowohl die noch fehlenden geodaetischen, als die in Italien erforderlichen astronomischen Operationen ausführen zu lassen. Diese Arbeiten machen den Inhalt des vorliegenden, wie es scheint von den Astronomen \textsc{Carlini} und \textsc{Plana} gemeinschaftlich redigirten Werks aus.

Der erste Theil ist ausschliesslich den geodaetischen Operationen gewidmet. Die beiden östlichen Endpunkte des Dreiecksnetzes in Frankreich, der Mont Colombier und der Mont Granier (unweit Chambery) bilden die Seite, von welcher die neue Messung ausgehen und bis zur westlichsten Seite des Netzes in der Lombardei, Massé --- Superga (bei Turin) fortgeführt werden musste. Man hätte erwarten sollen, dass in diesem Terrain, wo sich die höchsten Gebirge von Europa befinden, die Bildung grossartiger Dreiecke leicht, und eine sehr kleine Anzahl von Zwischenpunkten --- die Entfernung des Mont Granier von Superga beträgt nur 150000 Meter --- zur Verbindung hinreichend gewesen wäre. Allein gerade umgekehrt hatte man auf dieser mässigen Strecke mit den grössten Schwierigkeiten zu kämpfen, insofern die Spitzen der höheren Berge gar nicht oder schwer zugänglich sind, die Baumaterialien für die Signale nur mit grösster Anstrengung hinaufgeschafft werden können, und die heftigen Stürme sowohl diese Sig%
